%% kapa: exam=1 round=35 year=2024 period=1 form=A subject=economics subject_ko=경제학원론 seq=22 number=62
%% source: Q-Net 감정평가사 기출문제 (cst00309) artlSeq=5213229 file=raw/1008/5213229/1교시 문제지 원본.pdf
%% conversion: kordoc 4.24.51 -> tools/gen1.py (auto); review=proofread-sonnet; flags=none
\begin{question}{62}
2023년에 기업 A는 한국에서 생산한 부품 100억 달러를 베트남 현지 공장에 수출하였다. 같은 해에 베트남 현지 공장에서 그 부품을 조립하여 소비재 완제품 200억 달러를 만들어 그 중 50억 달러는 한국에 수출하고, 140억 달러는 미국에 수출하였으며 10억 달러는 재고로 남았다. 이월된 재고 10억 달러는 2024년 베트남 국내에서 모두 판매되었다. 이에 관한 설명으로 옳은 것은?
\begin{choices}
  \item 2023년 한국의 GDP는 50억 달러이다.
  \item 2023년 베트남의 GDP는 200억 달러이다.
  \item 2023년 베트남의 투자는 10억 달러이다.
  \item 2023년 베트남의 순수출은 190억 달러이다.
  \item 2024년 베트남의 소비와 GDP는 각각 10억 달러이다.
\end{choices}
\end{question}
\answer{3}
